\section{Introduction}
\label{sec:introduction}

Supervised learning, arguably the most popular branch of machine learning, involves estimating a mapping from data samples ($x$) to target variables ($y$). A common formulation of this task is the estimation of the conditional probability $p(y|x)$ from data through learning associations between $y$ and underlying factors of variation of $x$. However, data often contains nuisance factors ($z$) that are irrelevant to the prediction of $y$ from $x$ and estimation of $p(y|x)$ in such cases leads to overfitting when the model \emph{incorrectly} learns to associate some $z$ with $y$. Thus, when applied to new data containing unseen variations of $z$, trained models perform poorly. For example, a nuisance factor in the case of face recognition in images is the lighting condition the photograph was captured in, and a recognition model that associates lighting with subject identity is expected to perform poorly. Developing machine learning methods that are \emph{invariant} to nuisance factors has been a long-standing problem in machine learning; studied under various names such as ``feature selection'', ``robustness through data augmentation'' and ``invariance induction''.

While deep neural networks (DNNs) have outperformed traditional methods at highly sophisticated and challenging supervised learning tasks, providing better estimates of $p(y|x)$, they are prone to the same problem of incorrectly learning associations between $z$ and $y$. An architectural solution to this problem is the development of neural network units that capture specific forms of information, and thus are inherently invariant to certain nuisance factors~\cite{bib:drl,bib:sai}. For example, convolutional operations coupled with pooling strategies capture shift-invariant spatial information while recurrent operations robustly capture high-level trends in sequential data. However, this approach requires significant effort for \emph{engineering} custom network modules and layers to achieve invariance to \emph{specific} nuisance factors, making it inflexible~\cite{bib:sai}. A different but popularly adopted solution to the problem of nuisance factors is the use of data augmentation where synthetic versions of real data samples are generated, during training, with \emph{specific} forms of variation~\cite{bib:drl}. For example, rotation, translation and additive noise are typical methods of augmentation used in computer vision, especially for classification and detection tasks. However, models trained na\"ively on the augmented dataset become robust to \emph{limited} forms of nuisance by learning to associate every \emph{seen} variation of such factors to the target variables. Consequently, such models perform poorly when applied to data exhibiting \emph{unseen} nuisance variations, such as face images at previously unseen pose angles.% Moreover, data augmentation requires domain knowledge, which is not always readily available. Our approach is inherently robust to incomplete domain knowledge of nuisance factors, as well.

A related but more systematic solution to this problem is the approach of invariance-induction by guiding neural networks through specialized training mechanisms to discard \emph{known} nuisance factors from the learned latent representation of data that is used for prediction. Models trained in this fashion become robust by exclusion rather than inclusion and are, therefore, expected to perform well even on data containing variations of specific nuisance factors that were not seen during training. For example, a face recognition model trained explicitly to not associate lighting conditions with the identity of the person is expected to be more robust to lighting conditions than a similar model trained na\"ively on images of subjects under \emph{certain} different lighting conditions~\cite{bib:sai}. This research area has, therefore, garnered tremendous interest recently~\cite{bib:dann,bib:nnmmd,bib:vfae,bib:sai}. However, a shortcoming of this approach, is the requirement of domain knowledge of possible nuisance factors and their variations, which is often hard to find~\cite{bib:drl}. Additionally, this solution to invariance applies only to cases where annotated data is available for each nuisance factor, such as labeled information about the lighting condition of each image in the face recognition example, which is often not the case.

We present a novel unsupervised framework for invariance induction that overcomes the drawbacks of previous methods. Our framework promotes invariance through separating the underlying factors of variation of $x$ into two latent embeddings: $e_1$, which contains all the information required for predicting $y$, and $e_2$, which contains other information irrelevant to the prediction task. While $e_1$ is used for predicting $y$, a noisy version of $e_1$, denoted as $\tilde{e}_1$, and $e_2$ are used to reconstruct $x$. This creates a \textit{competitive} scenario where the reconstruction module tries to pull information into $e_2$ (because $\tilde{e}_1$ is unreliable) while the prediction module tries to pull information into $e_1$. The training objective is augmented with a disentanglement term that ensures that $e_1$ and $e_2$ do not contain redundant information. In our adversarial instantiation of this generalized framework, disentanglement is achieved between $e_1$ and $e_2$ in a novel way through two adversarial \textit{disentanglers} --- one that aims to predict $e_2$ from $e_1$ and another that does the inverse. The parameters of the combined model are learned through adversarial training between (a) the encoder, the predictor and the decoder, and (b) the disentanglers. The framework makes no assumptions about the data, so it can be applied to any prediction task without loss of generality, be it binary/multi-class classification or regression. Unlike existing methods, the proposed method does not require annotation of nuisance factors or specialized domain knowledge. We provide results on three tasks involving a diverse collection of datasets --- (1) invariance to inherent nuisance factors, (2) effective use of synthetic data augmentation for learning invariance and (3) domain adaptation. Our unsupervised framework outperforms existing approaches for invariance induction, which are supervised, on all of them.